%% kapa: exam=1 round=36 year=2025 period=2 form=A subject=accounting subject_ko=회계학 seq=20 number=60
%% source: Q-Net 감정평가사 기출문제 (cst00309) artlSeq=5231525 file=raw/1008/5231525/2025년 제36회 감정평가사 1차시험 2교시 문제지 원본.pdf
%% conversion: kordoc 4.24.51 -> tools/gen1.py (auto); review=proofread-sonnet; flags=none
\begin{question}{60}
20×1년 2월 초 영업을 개시한 (주)감평은 제품하자보증을 실시하고 있다. 제품매출액은 20×1년 ￦200,000, 20×2년 ￦250,000이고, (주)감평은 20×1년 매출액의 5 \%, 20×2년 매출액의 6 \%에 해당하는 무상수리비용이 발생할 것으로 추정하고 있다. 연도별 무상수리비용 실제 발생액은 다음과 같다.
\par{\small\centering
\begin{tabular}{|c|c|c|}\hline
구분 & 20×1년 & 20×2년 \\ \hline
20×1년 판매분 & ￦2,500 & ￦3,000 \\ \hline
20×2년 판매분 & - & 4,000 \\ \hline
\end{tabular}
\par}
(주)감평의 제품하자보증기간이 판매일로부터 2년인 경우, 20×2년 말 제품보증충당부채 잔액은? (단, 제품보증은 확신유형의 보증이고, 추정치는 최선의 추정치이며, 현재가치는 고려하지 않는다.)
\begin{choices}
  \item ￦4,500
  \item ￦6,500
  \item ￦9,500
  \item ￦11,000
  \item ￦15,500
\end{choices}
\end{question}
\answer{5}
